
\documentclass[utf8]{FrontiersinHarvard}
\usepackage{url,hyperref,lineno,microtype,subcaption}
\usepackage[onehalfspacing]{setspace}
\usepackage{amsmath}
\usepackage{booktabs}
\linenumbers
\def\keyFont{\fontsize{8}{11}\helveticabold }
\def\firstAuthorLast{Placeholder {et~al.}}
\def\Authors{A. Placeholder$^{1,*}$, B. Example$^{1}$ and C. Template$^{2}$}
\def\Address{$^{1}$Department of Verified Templates, Placeholder Institute, City, Country \\
$^{2}$Journal Systems Lab, Example University, City, Country}
\def\corrAuthor{A. Placeholder}
\def\corrEmail{verified@example.org}
\begin{document}
\onecolumn
\firstpage{1}
\title[Running Title]{Frontiers in Neurorobotics: Official Frontiers Submission Preview}
\author[\firstAuthorLast]{\Authors}
\address{}
\correspondance{}
\extraAuth{}
\maketitle
\begin{abstract}
\section{}
This submission-style preview is rendered with the imported official Frontiers template package and expanded into a realistic journal-manuscript skeleton. The content remains generic, but the page is intentionally arranged to show specialty metadata, correspondence, structured section flow, figure-table balance, and end statements in a form that is closer to a real Frontiers submission.
\tiny
\keyFont{ \section{Keywords:} neurorobotics, official template, submission preview, embodied intelligence, formatting check}
\end{abstract}
\section{Introduction}
This verified preview uses the official Frontiers class files from the publisher template package rather than a generic style engine. The opening section is written to resemble a concise introduction from a neurorobotics article so the reader can inspect how the title block, specialty metadata, abstract, and early body paragraphs relate on the first page.
\section{Related Work}
The related-work section is kept explicit so the preview resembles a true research article rather than a bare template sheet. This helps expose the vertical rhythm between literature framing, method description, and the first major figure or table.
\section{Methodology}
The methodology section introduces a representative equation and a figure placeholder so the official Frontiers package can be checked under a more realistic technical load.
\begin{equation}
J = \sum_{k=1}^{N} \left( \alpha \lVert e_k \rVert_2^2 + \beta \lVert u_k \rVert_2^2 \right).
\end{equation}
Here, $e_k$ denotes a placeholder tracking error, $u_k$ denotes an actuation term, and $\alpha,\beta > 0$ are weights included purely to make the preview read like a technical manuscript.
\begin{figure}[h]
\centering
\fbox{\parbox[c][2.2in][c]{0.82\linewidth}{\centering Frontiers Figure Placeholder\\[0.4em]\normalsize Experimental setup, embodiment diagram, or control pipeline}}
\caption{Representative figure block used to inspect the official Frontiers manuscript form, caption spacing, and the visual relationship between the methods text and the first major figure.}
\end{figure}
\section{Experiments}
The experiments section provides enough quantitative structure to make the preview more inspectable at a glance and closer to a practical submission manuscript.
\subsection{Experimental Setup}
This paragraph stands in for hardware conditions, datasets, simulation settings, and baseline models typically reported in a Frontiers submission.
\subsection{Quantitative Results}
\begin{table}[h]
\caption{Quantitative comparison table used to inspect the imported Frontiers package under a realistic neurorobotics-paper rhythm.}
\centering
\begin{tabular}{lcc}
\toprule
Method & Score & Latency \\
\midrule
Placeholder A & 0.91 & 14 ms \\
Placeholder B & 0.88 & 17 ms \\
Placeholder C & 0.84 & 20 ms \\
\bottomrule
\end{tabular}
\end{table}
\section{Conclusion and Discussion}
This page verifies the official Frontiers template path for this journal target while also making the section structure, figure flow, and reporting statements easier to inspect as a serious submission-style preview.
\section*{Data availability statement}
Placeholder data and code would be referenced here.
\section*{Author contributions}
All contributions are placeholders for template verification.
\section*{Conflict of interest}
The authors declare no conflict of interest in this verified preview.
\section*{Funding}
Placeholder funding statement.
\begin{thebibliography}{9}
\bibitem[Placeholder and Example(2026)]{ref1} Placeholder, A. and Example, B. (2026). Verified template note. \textit{Front. Neurorobot.} 1:1--2.
\end{thebibliography}
\end{document}
